\subsection{正测度的像}

设 X 是一个局部紧空间，\(\pi\) 是一个由 T 到 X 中的 \(\mu\) -可测映射。说偶 \((\pi,1)\) 是 \(\mu\) -适应的（§4, No.~1），等价于说对每个函数 \(f \in \mathcal{K}(X)\) ，函数 \(f \circ \pi\) 是本质 \(\mu\) -可积的。

命题 1. --- 设 \(\pi\) 是一个由 T 到局部紧空间 X 中的 \(\mu\) -可测映射。下列两个性质等价：

\begin{enumerate}
\def\labelenumi{\alph{enumi})}
\item
  对每个函数 \(f \in \mathcal{K}(\mathrm{X})\) ，\(f \circ \pi\) 是本质 \(\mu\) -可积的。
\item
  对每个紧集 \(\mathrm{K}\subset \mathrm{X}\) ，\(\overline{\pi} (\mathrm{K})\) 是本质 \(\mu\) -可积的。
\end{enumerate}

我们刚刚看到，a) 蕴涵偶 \((\pi,1)\) 是 \(\mu\) -适应的。于是（\(\S4\) , No.~4, 定理 2），对每个紧集 \(K\subset X\) ，函数 \(\varphi_{K}\circ\pi=\varphi_{A}\) （其中 \(A=\overline{\pi}(K)\) ）是本质 \(\mu\) -可积的，换言之，a) 蕴涵 b)。

反之，设 \(\overline{\pi}^{1}(K)\) 对 X 的每个紧子集 K 都是本质 \(\mu\) -可积的，我们来证明 a) 成立。实际上，设 S 是 f 的支集；由于 S 是紧的，置 \(A = \overline{\pi}^{1}(S)\) ，我们依假设有

\[
\int^ {\bullet} \left| f (\pi (t)) \right| d \mu (t) \leqslant \| f \| \int^ {\bullet} \varphi_ {\mathrm{S}} (\pi (t)) d \mu (t) = \| f \| \int^ {\bullet} \varphi_ {\mathrm{A}} (t) d \mu (t) <   + \infty .
\]

由于 \(f \circ \pi\) 是 \(\mu\) -可测的（Ch. IV, §5, No.~3, 定理 1），我们看出 \(f \circ \pi\) 是本质 \(\mu\) -可积的（§1, No.~3, 命题 9）。

性质 b) 显然等价于下列性质（因而后者也等价于 a)）：

\begin{enumerate}
\def\labelenumi{\alph{enumi})}
\setcounter{enumi}{2}
\tightlist
\item
  对每个点 \(x\) （\(X\) 中的），存在一个邻域 \(V\) （\(x\) 的），使得 \(\overline{\pi}^1(V)\) 是本质 \(\mu\) -可积的。
\end{enumerate}

定义 1. --- 设 \(\mu\) 是局部紧空间 T 上的一个正测度。称一个由 T 到局部紧空间 X 中的映射 \(\pi\) 为 \(\mu\) -逆紧的（或对测度 \(\mu\) 逆紧的），如果偶 \((\pi,1)\) 是 \(\mu\) -适应的，也就是说（§4, No.~1），如果 \(\pi\) 是 \(\mu\) -可测的且满足命题 1 的（等价的）条件。X 上的测度 \(\int \varepsilon_{\pi(t)} d\mu(t)\) 这时称为 \(\mu\) 在 \(\pi\) 下的像，记作 \(\pi(\mu)\) 。

于是若 \(\nu = \pi (\mu)\) ，则依定义，对 \(f\in \mathcal{K}(\mathbf{X})\) 有

\[
\int f (x) d \nu (x) = \int f (\pi (t)) d \mu (t).\tag{1}
\]

评注. --- 1) 若 \(\mu\) 是有界的（特别地，若 \(\mu\) 具有紧支集），则每个由 T 到 X 中的 \(\mu\) -可测映射都是 \(\mu\) -逆紧的（Ch. IV, §5, No.~3, 定理 1 及 No.~6, 定理 5）。

\begin{enumerate}
\def\labelenumi{\arabic{enumi})}
\setcounter{enumi}{1}
\item
  若 \(\pi\) 是 \(\mu\) -可测的，且对每个紧子集 \(K\) （\(X\) 的），\(\overline{\pi}^{1}(K)\) 是相对紧的，则 \(\pi\) 是 \(\mu\) -逆紧的（Ch. IV, §5, No.~5, 命题 7 及 No.~6, 定理 5）；特别地，每个由 \(T\) 到 \(X\) 中的逆紧连续映射（GT, I, §10, No.~2, 定理 1）都是 \(\mu\) -逆紧的，对每个正测度 \(\mu\) （\(T\) 上的）而言。更特别地，这对每个同胚 \(\pi\) （由 \(T\) 到 X 上的）都成立；此时测度 \(\nu = \pi(\mu)\) 就是 \(\mu\) 经 \(\pi\) 搬运到 X 上的那个测度（Ch. III, §1, No.~3）。
\item
  假设 X 的拓扑有一个可数基；则每个由 T 到 X 中的满足命题 1 的条件 b) 的映射 \(\pi\) 都是 \(\mu\) -可测的，从而是 \(\mu\) -逆紧的。只须应用 Ch. IV, §5, No.~5 的定理 4，并注意到此时 X 是可度量的（GT, IX, §2, No.~9, 命题 16 的推论），而且对 X 的拓扑相容的任何度量，每个闭球都是可数个紧集的并。
\end{enumerate}

